\chapter{Differentiating integrals, and derivatives from the definition}

\begin{orient}
\textbf{What this chapter is for.} Every rule you have met so far assumes that the function in front of you is an explicit formula, assembled from elementary pieces by arithmetic and composition. Two large families of problems supply no such formula, and in both of them the rules are not merely inconvenient but inapplicable. In the first family the function is defined by an integral: the variable sits in a limit, or inside the integrand, or in both places at once, and the antiderivative is typically unavailable and never needed. In the second family the function is assembled from pieces, or has a corner, or is assigned a value by fiat at the one point where its formula degenerates; there the difference quotient is the only instrument that answers the question. This chapter builds the full Leibniz integral rule from one multivariable chain rule so that its three terms stop being three things to memorise, then returns to the definition of the derivative and works through matching, one-sided limits, kinks, cusps, and the standard counterexamples that separate ``$f'(a)$ exists'' from ``$f'$ is continuous at $a$''. At the end you should be able to differentiate an integral-defined function on sight, and to recognise instantly when a problem is asking you to go back to first principles.
\end{orient}

\section{Functions defined by an integral}

Let $f$ be continuous and set $F(x)=\int_{a}^{x}f(t)\dd t$. The Fundamental Theorem says $F'(x)=f(x)$. Read that statement carefully, because its content is the engine of this entire section: integration of a continuous function produces a differentiable function, and differentiating undoes it \emph{pointwise}, with no reference whatever to a closed form for $F$. The function $\int_{0}^{x}\eu^{-t^{2}}\dd t$ has no elementary antiderivative, and its derivative is $\eu^{-x^{2}}$, known exactly and forever. That asymmetry between what is computable about $F$ and what is computable about $F'$ is precisely why integral-defined functions get set as derivative problems.

Everything in this section is that one theorem plus a chain rule. The order of presentation is deliberate: a variable upper limit, then two variable limits, then a variable inside the integrand as well. Each stage adds exactly one term, and the third stage is the general rule of which the first two are special cases. Learn the general rule and read the specialisations off it. Memorising three separate formulas is how the third term goes missing.

\tech{The Fundamental Theorem with a variable upper limit}{core}
\cue $F(x)=\int_{a}^{u(x)}f(t)\dd t$ with a constant lower limit and a non-trivial upper limit; typically an integrand with no elementary antiderivative, which is the setter's way of telling you not to integrate.
\method Compose the Fundamental Theorem with the chain rule. Write $G$ for any antiderivative of $f$, so that $F(x)=G(u(x))-G(a)$; then $F'(x)=G'(u(x))u'(x)$, that is,
\[
\dvop{x}\int_{a}^{u(x)}f(t)\dd t=f\bigl(u(x)\bigr)\,u'(x).
\]
The symbol $G$ is scaffolding: it appears in the justification and never in the calculation. If instead the \emph{lower} limit carries the variable, reverse the orientation first, $\int_{u(x)}^{a}f=-\int_{a}^{u(x)}f$, so the derivative is $-f(u(x))u'(x)$.

\begin{worked}
Take $F(x)=\displaystyle\int_{0}^{x^{3}}\eu^{-t^{2}}\dd t$, with $u(x)=x^{3}$ and $f(t)=\eu^{-t^{2}}$:
\[
F'(x)=\eu^{-(x^{3})^{2}}\cdot3x^{2}=3x^{2}\eu^{-x^{6}} .
\]
At $x=1.2$ the formula gives $0.21811592$, against a central difference of the numerically integrated $F$ of $0.21811592$. No attempt to evaluate $F$ was made, and none would have succeeded in closed form.

Now the variable in the lower limit. Let $H(x)=\displaystyle\int_{\sin x}^{2}\ln(1+t^{2})\dd t$. Flip the orientation first:
\[
H(x)=-\int_{2}^{\sin x}\ln(1+t^{2})\dd t,\qquad H'(x)=-\ln\bigl(1+\sin^{2}x\bigr)\cos x .
\]
At $x=\pi/4$ this is $-\ln\tfrac32\cdot\tfrac{\sqrt2}{2}=-0.2867071$, confirmed by central difference to seven digits.

Finally, a question that only looks like an integration problem. Where does $F(x)=\displaystyle\int_{0}^{x}(t^{2}-3t+2)\eu^{t}\dd t$ have local extrema? Do not integrate. The Fundamental Theorem gives $F'(x)=(x^{2}-3x+2)\eu^{x}=(x-1)(x-2)\eu^{x}$, and since $\eu^{x}>0$ the sign of $F'$ is the sign of $(x-1)(x-2)$: positive, negative, positive. So $F$ has a local maximum at $x=1$ and a local minimum at $x=2$, obtained without ever writing down $F$.
\end{worked}

\begin{trap}
Omitting $u'(x)$, which is the commonest error in the whole chapter, and mis-signing when the variable is downstairs: $\dvop{x}\int_{u(x)}^{a}f=-f(u(x))u'(x)$, with the minus coming from the orientation flip and not from anything about $f$. A third, subtler failure is trying to evaluate $F$ first. If the integrand has no elementary antiderivative, that attempt cannot terminate, and it was never required.
\end{trap}

\begin{seealsob}Both limits moving; the full Leibniz integral rule; outside-in chain-rule discipline.\end{seealsob}

Nothing changes structurally when the lower limit starts moving too. Split the integral at any convenient constant $c$ inside the interval, apply the previous technique to each half, and read off the answer. The reason the result is worth its own block is not the derivation but the design of the problems: setters choose the two limits so that the two contributions combine through an identity, and a solver who evaluates each piece numerically instead of looking for that identity misses the point of the question.

\tech{Both limits moving}{easy}
\cue Both limits carry $x$, very often reciprocally or as a power: $\int_{1/x}^{x}$, $\int_{x}^{x^{2}}$, $\int_{-x}^{x}$.
\method Split at a constant $c$ between them, $\int_{a(x)}^{b(x)}=\int_{c}^{b(x)}-\int_{c}^{a(x)}$, and differentiate each piece:
\[
\dvop{x}\int_{a(x)}^{b(x)}f(t)\dd t=f\bigl(b(x)\bigr)b'(x)-f\bigl(a(x)\bigr)a'(x).
\]
The upper limit enters with a plus and the lower limit with a minus, always. When the limits are reciprocal, expect a complementary identity such as $\arctan x+\arctan(1/x)=\pi/2$ for $x>0$ to collapse the two terms into one; look for it before reaching for a calculator.

\begin{worked}
Let $f(x)=\displaystyle\int_{1/x}^{x}\frac{\arctan t}{t}\dd t$ for $x>0$, and find $f'(3)$. Here $b(x)=x$, $b'=1$, $a(x)=1/x$, $a'=-1/x^{2}$, so
\[
f'(x)=\frac{\arctan x}{x}-\frac{\arctan(1/x)}{1/x}\cdot\left(-\frac{1}{x^{2}}\right)
=\frac{\arctan x}{x}+\frac{\arctan(1/x)}{x}=\frac{1}{x}\cdot\frac{\pi}{2}.
\]
So $f'(x)=\dfrac{\pi}{2x}$ for every $x>0$, and $f'(3)=\dfrac{\pi}{6}=0.5235988$, matching a central difference of the numerically integrated $f$ to seven digits. Note the three separate cancellations that made this clean: the $1/t$ in the integrand became $x$ when evaluated at $t=1/x$, that $x$ cancelled against the $1/x^{2}$ from $a'$, and the two arctangents summed to $\pi/2$. All three were engineered.

A second, where no identity is available and none is needed. With $K(x)=\displaystyle\int_{x}^{x^{2}}\ln t\dd t$ for $x>1$,
\[
K'(x)=\ln(x^{2})\cdot2x-\ln x\cdot1=4x\ln x-\ln x=(4x-1)\ln x ,
\]
so $K'(\eu)=4\eu-1=9.873127$, verified to seven digits.

A third, symmetric case, which is the one worth recognising on sight. For $G(x)=\displaystyle\int_{-x}^{x}f(t)\dd t$ the rule gives
\[
G'(x)=f(x)\cdot1-f(-x)\cdot(-1)=f(x)+f(-x),
\]
so an even integrand doubles and an odd integrand cancels to zero, recovering the standard symmetry facts as a one-line corollary. With $f(t)=t^{2}\eu^{t}$, which is neither even nor odd, $G'(x)=x^{2}(\eu^{x}+\eu^{-x})=2x^{2}\cosh x$, and $G'(1)=2\cosh1=3.0861613$, matching a central difference to seven digits.
\end{worked}

\begin{trap}
The lower-limit sign. In the first example the rule's minus and the negative $a'(x)=-1/x^{2}$ multiply to a plus, and the visible plus sign in the middle expression is exactly what makes a nervous solver ``correct'' it back to a minus. Write $b'$ and $a'$ down explicitly, with their signs, before substituting. The other error is splitting at a constant outside the interval of integration, which is legitimate for the algebra but can put you outside the domain of $f$.
\end{trap}

\begin{seealsob}The Fundamental Theorem with a variable upper limit; the full Leibniz integral rule; inverse-trigonometric complementary identities.\end{seealsob}

So far the variable has appeared only in the limits. The general case lets it appear in the integrand as well, and the resulting formula has three terms rather than two. It is not worth memorising as a formula. It is worth deriving once, because the derivation makes the three terms inevitable rather than arbitrary, and because seeing where each term comes from is what stops you dropping one.

Introduce a function of three independent slots,
\[
\Phi(x,y,z)=\int_{y}^{z}F(x,t)\dd t,
\]
and observe that the object we want is $\Phi\bigl(x,a(x),b(x)\bigr)$, a composition in which $x$ enters through three different doors. The multivariable chain rule therefore produces three terms, one per door. Through the third slot, $\partial\Phi/\partial z=F(x,z)$ by the Fundamental Theorem, multiplied by $b'(x)$. Through the second slot, $\partial\Phi/\partial y=-F(x,y)$, because $y$ is the lower limit, multiplied by $a'(x)$. Through the first slot the limits are frozen and only the integrand varies, giving $\int_{y}^{z}\partial F/\partial x\dd t$. Adding them yields the rule.

\begin{keynote}[Three doors, three terms]
The variable $x$ can enter an integral in exactly three ways: through the upper limit, through the lower limit, and through the integrand. Each entry contributes one term, and no term may be skipped merely because another one is present. If $x$ appears in all three places, your answer has three pieces; if in two, two; if in one, one. Counting the doors before differentiating is the checksum for this whole section.
\end{keynote}

The tree below is that count, drawn. It has no branch labelled ``look up a formula'', because there is only one formula and the branches merely record how many of its terms survive.

\begin{center}
\begin{tikzpicture}[node distance=6mm and 8mm]
\node[dtest] (a) {Does $x$ appear\\ inside the integrand?};
\node[dtest, below left=10mm and 2mm of a] (b) {Do both limits\\ move?};
\node[dtest, below right=10mm and 2mm of a] (c) {Is $x$ tied to $t$,\\ as $\lambda(x)t$?};
\node[dleaf, below=9mm of b] (d) {One or two boundary terms:\\ $f(b)b'-f(a)a'$};
\node[dleaf, below left=9mm and 0mm of c] (e) {Substitute $s=\lambda(x)t$;\\ carry the Jacobian};
\node[dleaf, below right=9mm and 0mm of c] (f) {Full Leibniz:\\ two boundary terms\\ plus $\int\pdv{F}{x}\dd t$};
\draw[dedge] (a) -- node[dlab,left]{no} (b);
\draw[dedge] (a) -- node[dlab,right]{yes} (c);
\draw[dedge] (b) -- node[dlab,left]{either way} (d);
\draw[dedge] (c) -- node[dlab,left]{yes} (e);
\draw[dedge] (c) -- node[dlab,right]{no} (f);
\end{tikzpicture}
\end{center}

\tech{The full Leibniz integral rule}{medium}
\cue Moving limits \emph{and} the variable present as a parameter inside the integrand: $\int_{x}^{x^{2}}\eu^{xt}/t\dd t$, $\int_{0}^{x}\sin(x-t)g(t)\dd t$.
\method All three contributions, every time:
\[
\dvop{x}\int_{a(x)}^{b(x)}F(x,t)\dd t
=F\bigl(x,b(x)\bigr)b'(x)-F\bigl(x,a(x)\bigr)a'(x)+\int_{a(x)}^{b(x)}\pdv{F}{x}(x,t)\dd t .
\]
The two boundary terms are evaluations, requiring no integration; the third is a fresh integral, usually easier than the original because differentiating in $x$ tends to simplify the $t$-dependence. Hypotheses: $F$ and $\partial F/\partial x$ continuous on a rectangle containing the swept region, and $a,b$ differentiable.

\begin{worked}
Let $f(x)=\displaystyle\int_{x}^{x^{2}}\frac{\eu^{xt}}{t}\dd t$ and find $f'(1)$. Identify the pieces: $F(x,t)=\eu^{xt}/t$, $b(x)=x^{2}$, $a(x)=x$, $\partial F/\partial x=\eu^{xt}$, the $t$ in the denominator having been consumed.

\emph{Upper boundary:} $F(x,x^{2})b'(x)=\dfrac{\eu^{x^{3}}}{x^{2}}\cdot2x=\dfrac{2\eu^{x^{3}}}{x}$.

\emph{Lower boundary:} $-F(x,x)a'(x)=-\dfrac{\eu^{x^{2}}}{x}$.

\emph{Interior:} $\displaystyle\int_{x}^{x^{2}}\eu^{xt}\dd t=\left[\frac{\eu^{xt}}{x}\right]_{t=x}^{t=x^{2}}=\frac{\eu^{x^{3}}-\eu^{x^{2}}}{x}$.

At $x=1$ the three pieces are $2\eu$, $-\eu$ and $0$, so $f'(1)=\eu=2.7182818$, confirmed by central difference to seven digits. The interior term vanishing at $x=1$ is a coincidence of the evaluation point, not a feature of the problem: at $x=1.2$ the three pieces are $9.3823064$, $-3.5172465$ and $1.1739067$, summing to $7.0389667$ against a central difference of $7.0389667$, and dropping the third would be a $16.7$ per cent error.

A case with fixed limits, where only the third door is open. Let $S(x)=\displaystyle\int_{0}^{1}\dfrac{\sin(xt)}{t}\dd t$. Both boundary terms are zero because $a'=b'=0$, and
\[
S'(x)=\int_{0}^{1}\cos(xt)\dd t=\frac{\sin x}{x}.
\]
At $x=1.7$ this is $0.5833322$, against a central difference of $0.5833322$. The original integrand has no elementary antiderivative in $t$; its $x$-derivative has an immediate one. That reversal is typical and is the reason the third term is usually the easy one.

Finally, the case that explains why this rule matters outside examination questions. Let
\[
y(x)=\int_{0}^{x}\sin(x-t)\,g(t)\dd t ,
\]
a convolution with $x$ in the upper limit \emph{and} in the integrand. The boundary term is $\sin(0)g(x)=0$ and the interior term is $\int_{0}^{x}\cos(x-t)g(t)\dd t$, so $y'(x)=\int_{0}^{x}\cos(x-t)g(t)\dd t$. Differentiating once more, the boundary term is now $\cos(0)g(x)=g(x)$ and the interior term is $-y(x)$, giving
\[
y''+y=g,\qquad y(0)=y'(0)=0 .
\]
The integral is the solution of a forced oscillator. Check it on $g(t)=t$: the integral evaluates to $y=x-\sin x$, so $y'=1-\cos x$ and $y''=\sin x$, and indeed $y''+y=x$. Numerically at $x=1.3$ the integral is $0.3364418$ and $x-\sin x=0.3364418$.
\end{worked}

\begin{trap}
Stopping after the two boundary terms because ``the parameter is already inside the integral'' and therefore feels accounted for. It is not: the boundary terms record how the \emph{region} moves, the interior term records how the \emph{integrand} changes, and these are independent effects. The second error is writing $\dv{F}{x}$ where the formula demands $\pdv{F}{x}$, and then differentiating $t$ as though it depended on $x$. Inside the third term, $t$ is a bound dummy and is held fixed.
\end{trap}

\begin{seealsob}Both limits moving; decoupling substitution before Leibniz; differentiating under the integral sign.\end{seealsob}

The three-term rule is general but it is not always the cheapest route. When the parameter and the dummy variable are algebraically tied together, so that the same expression in $x$ appears both in a limit and multiplying $t$, a substitution can move all of the $x$-dependence into the limits and a prefactor. What was a three-term Leibniz problem becomes a one-term Fundamental Theorem problem wrapped in a product or quotient rule, which is strictly less error-prone.

\tech{Decoupling substitution before Leibniz}{hard}
\cue The same function of $x$ appears in a limit \emph{and} inside the integrand, tied to $t$: $\int_{0}^{x^{2}}\ln(1+tx^{2})\dd t$, $\int_{0}^{x}g(t/x)\dd t$.
\method Substitute $s=\lambda(x)t$, where $\lambda$ is the offending factor. Then $\dd t=\dd s/\lambda(x)$, the prefactor $1/\lambda(x)$ comes outside, the integrand becomes a function of $s$ alone, and every trace of $x$ now lives in the upper limit and in that prefactor. Differentiate what remains with the product or quotient rule plus one Fundamental Theorem application. Use this whenever the substitution is available; fall back on the three-term rule when it is not.

\begin{worked}
Let $g(x)=\displaystyle\int_{0}^{x^{2}}\ln\bigl(1+tx^{2}\bigr)\dd t$ for $x>0$, and find $g'(1)$. The tie is the factor $x^{2}$ multiplying $t$, so put $s=tx^{2}$, giving $\dd t=\dd s/x^{2}$ and $t=x^{2}\Rightarrow s=x^{4}$:
\[
g(x)=\frac{1}{x^{2}}\int_{0}^{x^{4}}\ln(1+s)\dd s=\frac{(1+x^{4})\ln(1+x^{4})-x^{4}}{x^{2}},
\]
using $\int\ln(1+s)\dd s=(1+s)\ln(1+s)-s$. A numerical check before differentiating: at $x=1.3$ the original integral is $1.3895324$ and the closed form is $1.3895324$.

Now differentiate the closed form. Write $N(x)=(1+x^{4})\ln(1+x^{4})-x^{4}$, so $N'(x)=4x^{3}\ln(1+x^{4})+4x^{3}-4x^{3}=4x^{3}\ln(1+x^{4})$, a pleasant collapse. Then
\[
g'(x)=\frac{N'(x)x^{2}-2xN(x)}{x^{4}}=\frac{4x^{3}\ln(1+x^{4})}{x^{2}}-\frac{2N(x)}{x^{3}} .
\]
At $x=1$: $N(1)=2\ln2-1$ and $N'(1)=4\ln2$, so
\[
g'(1)=4\ln2-2(2\ln2-1)=4\ln2-4\ln2+2=2 .
\]
Verified: $2.0000000$. The logarithms cancel exactly, which is the signature of a designed problem; had you run the three-term rule you would have produced $\ln(1+x^{4})\cdot2x$, $0$, and $\int_{0}^{x^{2}}\frac{2xt}{1+tx^{2}}\dd t$, and the same cancellation would have happened one stage later and less visibly.

The extreme case, where decoupling removes the calculus entirely. Let $\Psi(x)=\displaystyle\int_{0}^{x}\sin\frac{t}{x}\dd t$ for $x>0$. Put $u=t/x$, so $\dd t=x\dd u$ and the limits become $0$ and $1$:
\[
\Psi(x)=x\int_{0}^{1}\sin u\dd u=x\bigl(1-\cos1\bigr).
\]
The function is linear, and $\Psi'(x)=1-\cos1=0.4596977$ for every $x$, confirmed numerically at $x=2.7$ where the integral itself is $1.2411838=2.7(1-\cos1)$. The three-term rule would have produced $\sin1$ from the boundary and $-\frac{1}{x}\int_{0}^{x}\frac{t}{x}\cos\frac tx\dd t$ from the interior, which is correct and takes a further integration by parts to reach the same constant. Any integrand of the form $g(t/x)$ decouples this way.
\end{worked}

\begin{trap}
Forgetting the Jacobian. The factor $1/\lambda(x)$ is not decoration: it depends on $x$ and is exactly what forces the quotient rule at the end. Dropping it produces an answer that is wrong by a smooth factor, which no sanity check on magnitude will catch. The second trap is substituting when $\lambda(x)$ can vanish on the domain, which destroys the change of variable at that point.
\end{trap}

\begin{seealsob}The full Leibniz integral rule; algebraic normalisation before differentiating; differentiating under the integral sign.\end{seealsob}

There is a second reason to differentiate an integral, and it inverts the logic of everything above. Instead of being handed a function defined by an integral and asked for its derivative, you are handed a definite integral depending on a parameter, and differentiating in that parameter is the fastest way to \emph{evaluate} it. The limits are fixed, so only the third door is open, and the rule reduces to a single interior term. This is the technique usually named after Feynman, though it is Leibniz's rule with two terms switched off.

\tech{Differentiating under the integral sign}{hard}
\cue A definite integral over fixed limits, often infinite, whose integrand carries a parameter $a$; and a question about $\dv{}{a}$ of it, or an evaluation that resists every substitution.
\method With fixed limits the boundary terms vanish and
\[
\dvop{a}\int_{a_{0}}^{b_{0}}F(a,t)\dd t=\int_{a_{0}}^{b_{0}}\pdv{F}{a}(a,t)\dd t .
\]
Two uses. Forwards: differentiate a known parametric identity to get a new one for free. Backwards: choose the parameter so that $\partial_{a}F$ is elementary, evaluate the easy integral to get $I'(a)$, then integrate in $a$ and fix the constant from a value of $a$ where $I$ is obvious. On an infinite interval the interchange needs justification: uniform convergence of $\int\partial_{a}F$ on a neighbourhood of $a$, or a dominating integrable function.

\begin{worked}
Forwards, on a standard identity. For $a\geq0$,
\[
\int_{0}^{\infty}\frac{\ln(1+a^{2}x^{2})}{1+x^{2}}\dd x=\pi\ln(1+a),
\]
which at $a=2$ evaluates to $\pi\ln3=3.4513923$ against a numerical value of $3.4513$. Differentiating the closed form gives $\dfrac{\pi}{1+a}$, so the $a$-derivative at $a=2$ is $\dfrac{\pi}{3}=1.0471976$. Nobody should differentiate the integrand here: the closed form is already available and the derivative of the answer is one line, whereas $\int_{0}^{\infty}\frac{2ax^{2}}{(1+a^{2}x^{2})(1+x^{2})}\dd x$ needs partial fractions to reach the same place.

Backwards, where the closed form is what we are after. Let $J(a)=\displaystyle\int_{0}^{\infty}\eu^{-ax}\frac{\sin x}{x}\dd x$ for $a>0$. Differentiating under the sign kills the awkward $1/x$:
\[
J'(a)=-\int_{0}^{\infty}\eu^{-ax}\sin x\dd x=-\frac{1}{1+a^{2}} ,
\]
so $J(a)=-\arctan a+C$. As $a\to\infty$ the integrand is crushed and $J\to0$, forcing $C=\pi/2$, hence $J(a)=\arctan(1/a)$. At $a=1.5$ the formula gives $0.5880026$ and the numerical integral gives $0.5880026$; the derivative there is $-1/(1+2.25)=-0.3076923$. Letting $a\to0^{+}$ recovers $\int_{0}^{\infty}\frac{\sin x}{x}\dd x=\pi/2$ as a bonus.

The moment generator, which is the technique's most reused instance. From
\[
I(a)=\int_{0}^{\infty}\eu^{-ax^{2}}\dd x=\frac12\sqrt{\frac\pi a}\qquad(a>0),
\]
differentiating in $a$ brings down a factor $-x^{2}$ on the left and differentiates $\tfrac12\sqrt\pi\,a^{-1/2}$ on the right:
\[
\int_{0}^{\infty}x^{2}\eu^{-ax^{2}}\dd x=\frac14\sqrt{\frac{\pi}{a^{3}}} .
\]
At $a=1$ that is $\dfrac{\sqrt\pi}{4}=0.4431135$, against a numerical value of $0.4431135$. Differentiating again gives $\int_{0}^{\infty}x^{4}\eu^{-ax^{2}}\dd x=\tfrac38\sqrt{\pi/a^{5}}$, which at $a=1$ is $0.6646702$, again confirmed numerically. Every even moment of the Gaussian falls out of one identity and repeated differentiation, and the parameter $a$ was introduced solely to have something to differentiate with respect to.
\end{worked}

\begin{trap}
Swapping the derivative and the integral on an improper integral without a word of justification. The interchange can genuinely fail, and the standard counterexamples have an integrand whose mass escapes to infinity as the parameter moves. The other, more practical, trap is answering the wrong question: if the problem asks only for $I'(a)$, evaluating $I(a)$ first is wasted labour, and the derivative is often far simpler than the function.
\end{trap}

\begin{seealsob}The full Leibniz integral rule; special functions defined by a differential equation; partial derivatives and mixed partials.\end{seealsob}

The last member of this family turns the machinery around once more. The unknown is now the function inside the integral, and the integral relation itself is the equation to be solved. Differentiation is the tool that converts an integral equation into something local, and the classical case where the naive differentiation fails is worth seeing, because it shows what the hypotheses of the Fundamental Theorem are actually protecting.

\tech{A function defined implicitly by an integral equation}{hard}
\cue The unknown $y$ appears only inside an integral with a variable limit: $\int_{0}^{x}K(x-t)y(t)\dd t=h(x)$, a Volterra convolution equation. The Abel kernel $K(u)=u^{-1/2}$ is the classic case.
\method If $K$ is continuous, differentiate the whole equation in $x$ by the Leibniz rule; the boundary term is $K(0)y(x)$ and the interior term is $\int_{0}^{x}K'(x-t)y(t)\dd t$, converting the integral equation into a simpler one or into an ODE. If $K$ is singular at $u=0$, as the Abel kernel is, that step is illegal and you use the Laplace transform instead: convolution becomes multiplication, so $\widehat K(s)\widehat y(s)=\widehat h(s)$, and with $\widehat{u^{-1/2}}(s)=\sqrt{\pi/s}$ the equation inverts in one line.

\begin{worked}
Solve $\displaystyle\int_{0}^{x}\frac{y(t)}{\sqrt{x-t}}\dd t=x^{2}$ for $x>0$, and evaluate $\pi y'(4)$.

The kernel is singular at $t=x$, so transform. With $\widehat h(s)=2/s^{3}$ and $\widehat K(s)=\sqrt{\pi/s}$,
\[
\widehat y(s)=\frac{2}{s^{3}}\sqrt{\frac{s}{\pi}}=\frac{2}{\sqrt\pi}\,s^{-5/2}.
\]
Inverting with $s^{-\nu}\mapsto x^{\nu-1}/\Gfun(\nu)$ and $\Gfun(5/2)=\tfrac34\sqrt\pi$,
\[
y(x)=\frac{2}{\sqrt\pi}\cdot\frac{x^{3/2}}{\tfrac34\sqrt\pi}=\frac{8x^{3/2}}{3\pi}.
\]
Check it directly rather than trusting the transform: substituting $t=x-u^{2}$ turns the left side into $2\int_{0}^{\sqrt x}y(x-u^{2})\dd u$, which at $x=2$ evaluates numerically to $4.0000000$, exactly $x^{2}$. Then $y'(x)=\dfrac{8}{3\pi}\cdot\dfrac32 x^{1/2}=\dfrac{4\sqrt x}{\pi}$, so $y'(4)=\dfrac{8}{\pi}$ and
\[
\pi y'(4)=8 .
\]

Contrast the regular case, where differentiation is legal and finishes the job on the spot. Solve $\displaystyle\int_{0}^{x}\eu^{x-t}y(t)\dd t=x$. Here $K(u)=\eu^{u}$ is continuous, so apply the Leibniz rule to both sides: the boundary term is $K(0)y(x)=y(x)$, the interior term is $\int_{0}^{x}\eu^{x-t}y(t)\dd t$, which the original equation says equals $x$, and the right side differentiates to $1$. Hence
\[
y(x)+x=1,\qquad y(x)=1-x ,
\]
with no transform and no series. Substituting back, $\int_{0}^{x}\eu^{x-t}(1-t)\dd t=\eu^{x}\cdot x\eu^{-x}=x$, which at $x=1.7$ evaluates numerically to $1.7000000$. The lesson is the dividing line: a continuous kernel yields to differentiation, a kernel singular at $u=0$ does not.
\end{worked}

\begin{trap}
Differentiating under the sign when the kernel blows up at the endpoint. For the Abel equation the interior term $\int_{0}^{x}\partial_{x}(x-t)^{-1/2}y(t)\dd t$ diverges and the boundary term $K(0)y(x)$ is infinite, so the manipulation produces nonsense that looks formally correct. Either smooth first by integrating once, or transform. The second trap is forgetting that a Volterra equation carries its own initial condition, so no constant of integration is free to be chosen later.
\end{trap}

\begin{seealsob}The full Leibniz integral rule; differentiating under the integral sign; the Fundamental Theorem with a variable upper limit.\end{seealsob}

\subsection{Inverses of integral-defined functions}

A function defined by an integral of a positive integrand is strictly increasing, hence invertible, and questions about the inverse are then answerable even though neither the function nor its inverse can be written down. The tool is the inverse-function rule, $\bigl(G^{-1}\bigr)'(y)=\dfrac{1}{G'(x)}$ evaluated at the $x$ with $G(x)=y$, and the Fundamental Theorem supplies $G'$ for free.

Let $G(x)=\displaystyle\int_{0}^{x}\eu^{t^{2}}\dd t$. The integrand is positive, so $G$ is strictly increasing on $\R$ and $G^{-1}$ exists. Since $G(0)=0$,
\[
\bigl(G^{-1}\bigr)'(0)=\frac{1}{G'(0)}=\frac{1}{\eu^{0}}=1 .
\]
At a less convenient point the only extra work is locating $x$. Numerically $G(1)=1.4626517$, so
\[
\bigl(G^{-1}\bigr)'(1.4626517)=\frac{1}{\eu^{1}}=0.3678794 .
\]
The second derivative follows from the standard formula $\bigl(G^{-1}\bigr)''=-G''/(G')^{3}$, and here $G''(x)=2x\eu^{x^{2}}$, which vanishes at $x=0$; so $G^{-1}$ has an inflection at $y=0$, a fact about a function nobody can write down.

The same reading explains a familiar object. With $L(x)=\int_{1}^{x}\dd t/t$ the Fundamental Theorem gives $L'(x)=1/x$, so $\bigl(L^{-1}\bigr)'(y)=1/L'(x)=x=L^{-1}(y)$: the inverse is its own derivative, which is the definition of the exponential. The logarithm defined as an integral and the exponential defined as its inverse are constructed exactly this way.

\subsection{Second derivatives, and shape without a formula}

Nothing stops you differentiating twice. If $F(x)=\int_{a}^{x}f(t)\dd t$ then $F'=f$ and $F''=f'$, so the entire apparatus of the next chapter, monotonicity from the sign of the first derivative and concavity from the sign of the second, applies to a function you cannot write down. This is worth isolating because it is how integral-defined functions appear in shape questions, and because the answers are usually cleaner than the function.

Take $F(x)=\displaystyle\int_{0}^{x}(t^{2}-3t+2)\eu^{t}\dd t$ again. We already have $F'(x)=(x-1)(x-2)\eu^{x}$, so $F$ rises on $(-\infty,1)$, falls on $(1,2)$, and rises thereafter. For concavity, differentiate again by the product rule:
\[
F''(x)=(2x-3)\eu^{x}+(x^{2}-3x+2)\eu^{x}=(x^{2}-x-1)\eu^{x},
\]
which vanishes at $x=\dfrac{1\pm\sqrt5}{2}$, that is at $1.6180340$ and $-0.6180340$. Both are simple roots of the quadratic, so the sign genuinely changes and both are inflection points. Note that neither coincides with a critical point, and that the golden ratio has appeared out of an integrand chosen for its ordinariness.

When the upper limit is a function of $x$, the second derivative needs a product rule as well as a chain rule. From $F(x)=\int_{a}^{u(x)}f$, differentiating $F'=f(u)u'$ gives
\[
F''(x)=f'\bigl(u(x)\bigr)u'(x)^{2}+f\bigl(u(x)\bigr)u''(x),
\]
which is Faa di Bruno's formula at order two and is the shape you should expect: a term with the inner derivative squared, and a term with the inner second derivative. For $F(x)=\int_{0}^{x^{3}}\eu^{-t^{2}}\dd t$ this reads $F''(x)=-2x^{3}\eu^{-x^{6}}\cdot9x^{4}+\eu^{-x^{6}}\cdot6x$, that is $6x\eu^{-x^{6}}\bigl(1-3x^{6}\bigr)$, so the graph has inflections at $x=0$ and at $x=\pm3^{-1/6}$.

\section{Back to the definition}

The rest of this chapter is about the other kind of failure. In the integral family the rules did not apply because the function had no formula; here the function has a formula, or several, but the formula is not differentiable-by-rule at the point in question. The remedy is always the same and it is always available:
\[
f'(a)=\lim_{h\to0}\frac{f(a+h)-f(a)}{h},
\]
computed as an honest limit. This is not a fallback for the unlucky. It is the definition, and there are questions that only it can answer.

Four situations force you back to it. The function is defined by cases and $a$ is a junction point. The function is defined by a rule that degenerates at $a$, with $f(a)$ supplied separately. The function involves absolute values, floors, or fractional powers whose derivative formula fails at $a$. Or the function is pathological on purpose, as in the counterexamples that distinguish differentiability from continuous differentiability. The first technique below covers the raw computation; the rest specialise it.

\tech{The derivative from first principles}{medium}
\cue The problem says ``from the definition''; or the function is defined by cases, by conditions on rationality, or by a limit; or the rules produce an expression that is undefined at the very point requested.
\method Form the difference quotient and evaluate the limit honestly. The workhorse supporting limits are $\lim_{h\to0}\frac{\sin h}{h}=1$, $\lim_{h\to0}\frac{\eu^{h}-1}{h}=1$, $\lim_{h\to0}\frac{\ln(1+h)}{h}=1$, and $\lim_{h\to0}\frac{(1+h)^{r}-1}{h}=r$. When the numerator cannot be factored, squeeze: if $\abs{f(a+h)-f(a)}\leq C\abs{h}^{1+\epsilon}$ for some $\epsilon>0$, then $f'(a)=0$ immediately, with no algebra at all.

\begin{worked}
The squeeze in its purest form. Let $D$ be the Dirichlet function, equal to $1$ at rational points and $0$ elsewhere, and set $f(x)=x^{2}D(x)$. This function is discontinuous at every point except $0$, so no rule of any kind applies anywhere. At the origin, however, $f(0)=0$ and
\[
\left|\frac{f(h)-f(0)}{h}\right|=\left|\frac{h^{2}D(h)}{h}\right|\leq\abs{h}\To0 ,
\]
so $f'(0)=0$ exists. A function can be differentiable at exactly one point.

A case where the two sides disagree, so the answer is that the derivative does not exist. Let $f(x)=\dfrac{x}{1+\eu^{1/x}}$ for $x\neq0$ with $f(0)=0$. The quotient is
\[
\frac{f(h)-f(0)}{h}=\frac{1}{1+\eu^{1/h}} .
\]
As $h\to0^{+}$, $1/h\to+\infty$ and $\eu^{1/h}\to\infty$, so the quotient tends to $0$. As $h\to0^{-}$, $1/h\to-\infty$ and $\eu^{1/h}\to0$, so the quotient tends to $1$. Both one-sided derivatives exist and are finite, $f'_{+}(0)=0$ and $f'_{-}(0)=1$, and they differ: $f$ has a corner at the origin. Observe that $f$ \emph{is} continuous there, since $\abs{f(x)}\leq\abs{x}$; continuity was never the issue.

And a case where the algebra has to be done rather than squeezed. Let $f(x)=\abs{x}^{3}$ and compute $f'(0)$ and $f''(0)$ from the definition. First,
\[
\frac{f(h)-f(0)}{h}=\frac{\abs{h}^{3}}{h}=h\abs{h}\To0 ,
\]
so $f'(0)=0$. Away from $0$ the rules give $f'(x)=3x\abs{x}$, which is continuous at $0$ with value $0$, so $f$ is $C^{1}$. Repeating,
\[
\frac{f'(h)-f'(0)}{h}=\frac{3h\abs{h}}{h}=3\abs{h}\To0 ,
\]
so $f''(0)=0$ and $f''(x)=6\abs{x}$, again continuous. But $f'''(x)=6\sgn x$ has a genuine jump at $0$, so $f$ is $C^{2}$ and no better. The bars did not obstruct anything until the third derivative.
\end{worked}

\begin{trap}
Assuming that a function built from differentiable pieces must be differentiable at the junction, and evaluating a piece's derivative formula at the junction instead of taking a one-sided limit of the difference quotient. The second trap is cancelling $h$ out of the numerator before checking that a factor of $h$ is genuinely there; if the numerator is $O(h^{1/2})$, the quotient blows up and the derivative does not exist.
\end{trap}

\begin{seealsob}Recognising a limit as a derivative in disguise; one-sided derivatives, corners and cusps; piecewise matching for continuity and differentiability.\end{seealsob}

The definition can also be run backwards, and doing so is one of the highest-value recognitions in elementary calculus. A limit that looks intractable is very often a difference quotient wearing a disguise, and once it is recognised the answer is a single derivative evaluation. The move costs nothing and it works in cases where L'Hopital's rule is not licensed, since the definition needs $g$ differentiable only \emph{at} the point, while L'Hopital needs $g'$ to exist in a whole punctured neighbourhood.

\tech{Recognising a limit as a derivative in disguise}{easy}
\cue A limit of shape $\lim_{x\to a}\dfrac{g(x)-g(a)}{x-a}$, or $\lim_{h\to0}\dfrac{g(a+\alpha h)-g(a+\beta h)}{h}$, or $\lim_{n\to\infty}n\bigl(g(a+c/n)-g(a)\bigr)$; typically an indeterminate $0/0$ in which one specific number keeps reappearing as an argument.
\method Match the shape to the definition and identify $g$ and $a$. The symmetric two-point version follows by adding and subtracting $g(a)$:
\[
\lim_{h\to0}\frac{g(a+\alpha h)-g(a+\beta h)}{h}=(\alpha-\beta)\,g'(a).
\]
The $n\to\infty$ version is the same statement with $h=1/n$. Read it in whichever direction is cheaper: forwards it evaluates a hard limit by one differentiation, backwards it computes a derivative that the rules cannot reach.

\begin{worked}
Evaluate $\displaystyle\lim_{h\to0}\frac{\sin\bigl(\tfrac{\pi}{6}+3h\bigr)-\sin\bigl(\tfrac{\pi}{6}-h\bigr)}{h}$. The repeated argument is $\pi/6$, so $g=\sin$, $a=\pi/6$, $\alpha=3$, $\beta=-1$, and the limit is
\[
(3-(-1))\cos\frac{\pi}{6}=4\cdot\frac{\sqrt3}{2}=2\sqrt3=3.4641016 ,
\]
against a numerical quotient at $h=10^{-6}$ of $3.4640996$. Expanding both sines with the addition formula reaches the same answer in six more lines.

A sequence version. Find $\displaystyle\lim_{n\to\infty}n\bigl(2^{1/n}-1\bigr)$. Write $2^{1/n}=g(1/n)$ with $g(u)=2^{u}$ and note $g(0)=1$, so the expression is $\dfrac{g(1/n)-g(0)}{1/n}\to g'(0)=\ln2=0.6931472$. Numerically at $n=10^{7}$ the value is $0.6931472$.

A version with a tangent, where the coefficient does the work. Consider $\displaystyle\lim_{h\to0}\frac{\tan\bigl(\tfrac\pi4+h\bigr)-\tan\bigl(\tfrac\pi4-h\bigr)}{h}$. The repeated argument is $\pi/4$, so $g=\tan$, $\alpha=1$, $\beta=-1$, and the limit is $2\sec^{2}(\pi/4)=2\cdot2=4$, against a numerical quotient of $4.0000000$. Expanding both tangents with the addition formula produces a ratio of quadratics in $\tan h$ and takes half a page.

And backwards, to compute something the rules handle badly. What is $\dvop{x}x^{x}$ at $x=3$? The rules do reach this by logarithmic differentiation, but the definition reaches it too:
\[
\lim_{h\to0}\frac{(3+h)^{3+h}-3^{3}}{h}=3^{3}(1+\ln3)=56.662532 ,
\]
matching a symmetric difference of $56.662532$. The point is not that the definition is faster here, but that it is the same object: a derivative and this limit are two names for one number, and a problem may hand you either.
\end{worked}

\begin{trap}
Reaching for L'Hopital's rule reflexively. It is heavier than the definition and it has a hypothesis the definition does not: $g'$ must exist near $a$, not just at $a$. On a function like $x^{2}D(x)$ at the origin, L'Hopital is simply unavailable while the definition answers instantly. The second trap is misreading $\alpha$ and $\beta$: the coefficient is $\alpha-\beta$, so a minus sign inside the second argument becomes a plus in the answer.
\end{trap}

\begin{seealsob}The derivative from first principles; the derivative at a point of continuous extension; linearisation and differentials.\end{seealsob}

A harder variant arises when the formula for $f$ is not merely awkward at $a$ but indeterminate there, with $f(a)$ assigned by continuous extension. Now the difference quotient is a $0/0$ limit whose numerator is itself the difference of an indeterminate form from its limiting value, and the only reliable instrument is a series expansion carried far enough.

\tech{The derivative at a point of continuous extension}{hard}
\cue $f$ has a removable indeterminacy at exactly the point where $f'$ is requested, usually $1^{\infty}$ or $\infty^{0}$, with $f(a)$ supplied by fiat to fill the hole.
\method Do not differentiate the formula. Expand. Take logarithms and produce the Taylor expansion of $\ln f$ about $a$ to first order in $(x-a)$: if
\[
\ln f(x)=c_{0}+c_{1}(x-a)+O\bigl((x-a)^{2}\bigr),\quad\text{then}\quad f(a)=\eu^{c_{0}},\qquad f'(a)=f(a)\,c_{1}.
\]
The bookkeeping point is that an exponent of the form $1/(x-a)$ divides the expansion by $(x-a)$, so to obtain $c_{1}$ you must expand the bracket to \emph{second} order. Truncating at first order loses exactly the term you need.

\begin{worked}
Let $f(x)=\left(\dfrac{\ln x}{x-1}\right)^{1/(x-1)}$ for $x>0$, $x\neq1$, with $f(1)=\eu^{-1/2}$. Find $f'(1)$.

Set $u=x-1$. Then $\ln x=\ln(1+u)=u-\dfrac{u^{2}}{2}+\dfrac{u^{3}}{3}-\cdots$, so
\[
\frac{\ln x}{x-1}=1-\frac{u}{2}+\frac{u^{2}}{3}-\cdots=1+w,\qquad w=-\frac{u}{2}+\frac{u^{2}}{3}+O(u^{3}).
\]
Now $\ln(1+w)=w-\dfrac{w^{2}}{2}+O(w^{3})$, and $w^{2}=\dfrac{u^{2}}{4}+O(u^{3})$, so
\[
\ln\frac{\ln x}{x-1}=-\frac{u}{2}+\frac{u^{2}}{3}-\frac{u^{2}}{8}+O(u^{3})=-\frac{u}{2}+\frac{5u^{2}}{24}+O(u^{3}).
\]
Dividing by $u$ gives $\ln f=-\dfrac12+\dfrac{5u}{24}+O(u^{2})$, so $c_{0}=-\tfrac12$, confirming $f(1)=\eu^{-1/2}$, and $c_{1}=\tfrac{5}{24}$. Hence
\[
f'(1)=\eu^{-1/2}\cdot\frac{5}{24}=0.1263606 ,
\]
against a central difference of $0.1263606$. Note where the $5/24$ came from: $\tfrac13$ from the cubic term of $\ln(1+u)$ and $-\tfrac18$ from the quadratic correction in $\ln(1+w)$. Stop either expansion one order early and you get $\tfrac13$ or $-\tfrac18$ alone, both wrong.

A shorter instance of the same machine. Let $f(x)=(1+x)^{1/x}$ for $x>-1$, $x\neq0$, with $f(0)=\eu$. Then
\[
\ln f=\frac{\ln(1+x)}{x}=\frac1x\left(x-\frac{x^{2}}{2}+\frac{x^{3}}{3}-\cdots\right)=1-\frac{x}{2}+\frac{x^{2}}{3}-\cdots ,
\]
so $c_{0}=1$ and $c_{1}=-\tfrac12$, giving $f(0)=\eu$ as stated and
\[
f'(0)=\eu\cdot\left(-\frac12\right)=-\frac{\eu}{2}=-1.3591409 ,
\]
against a central difference of $-1.3591409$. Here the second-order term of $\ln(1+x)$ is what produced $c_{1}$, again one order deeper than a careless expansion would go.
\end{worked}

\begin{trap}
Expanding to first order only. The division by $(x-a)$ demotes every order by one, so first-order accuracy in $f'$ demands second-order accuracy in the bracket. The same arithmetic explains why applying L'Hopital once to $\frac{f(x)-f(a)}{x-a}$ typically stalls: the rule must be applied twice, and after the first application the expression is worse, not better.
\end{trap}

\begin{seealsob}The derivative from first principles; Maclaurin coefficients by series manipulation; recognising a limit as a derivative in disguise.\end{seealsob}

\section{Matching, corners, and the failure of the rules}

The three techniques of the previous section all computed a derivative that existed. The remaining question is when one exists at all, and the answer is always decided by comparing two one-sided limits. That comparison is the organising idea for everything below: matching problems choose parameters to force the two sides to agree, corner and cusp problems ask you to observe that they do not, and the final counterexample family shows that agreement of the one-sided derivatives is a strictly weaker condition than continuity of $f'$.

The most common examination form of the definition is not exotic at all. A function is given by two formulas meeting at a point, with unknown constants, and the constants must be chosen so that the join is smooth. This is two equations in two unknowns, and the only subtlety is which two equations.

\tech{Piecewise matching for continuity and differentiability}{core}
\cue $f$ given by two formulas meeting at $x=c$, with parameters $a,b$ to be determined so that $f$ is differentiable (or twice differentiable) at $c$.
\method Impose the conditions in order, because the second presupposes the first. \emph{Continuity:} $\lim_{x\to c^{-}}f(x)=\lim_{x\to c^{+}}f(x)=f(c)$. \emph{Differentiability:} $f'_{-}(c)=f'_{+}(c)$. Solve the resulting $2\times2$ system. The shortcut $\lim_{x\to c^{-}}f'(x)=\lim_{x\to c^{+}}f'(x)$ is legitimate only when $f$ is already continuous at $c$ and both one-sided limits of $f'$ exist; that is the usual situation with polynomial and exponential pieces, and it is what makes the method quick. When either hypothesis fails, go back to one-sided difference quotients.

\begin{worked}
The archetype first. Let $f(x)=x^{2}$ for $x\leq1$ and $f(x)=ax+b$ for $x>1$. Continuity at $1$: $1=a+b$. Differentiability: the left derivative is $2x$ at $x=1$, namely $2$, and the right derivative is $a$, so $a=2$ and therefore $b=-1$. The smooth continuation of $x^{2}$ by a line is its own tangent, $y=2x-1$, which is the geometric content of the answer.

A version where the parameters sit on the other side and the second derivative is also matched. Let
\[
f(x)=\begin{cases} ax^{2}+bx+1, & x\leq0,\\[2pt] \eu^{2x}\cos3x, & x>0.\end{cases}
\]
Continuity: the left limit is $1$ and the right limit is $\eu^{0}\cos0=1$, so continuity holds for every $a$ and $b$; the setter has spent that equation. Differentiability: the right piece has $y'=\eu^{2x}(2\cos3x-3\sin3x)$, so $f'_{+}(0)=2$, while $f'_{-}(0)=b$. Hence $b=2$. That exhausts the first-derivative condition and leaves $a$ free, so the question must be asking for $C^{2}$. Differentiating again,
\[
y''=\eu^{2x}\bigl(4\cos3x-6\sin3x\bigr)+\eu^{2x}\bigl(-6\sin3x-9\cos3x\bigr)=\eu^{2x}\bigl(-5\cos3x-12\sin3x\bigr),
\]
so $f''_{+}(0)=-5$, while $f''_{-}(0)=2a$. Therefore $a=-\tfrac52$, $b=2$. Numerically, the right piece has $y'(0)=2.0000000$ and $y''(0)=-5.0000001$ by central differences. The general pattern: matching $k$ derivatives consumes $k+1$ equations, counting continuity as the zeroth.

And a case where the derivative shortcut is illegal and must be replaced. Let $f(x)=x^{2}\sin(1/x)$ for $x<0$ and $f(x)=ax+b$ for $x\geq0$. Continuity at $0$: the left limit is $0$ by the squeeze $\abs{f}\leq x^{2}$, and the right value is $b$, so $b=0$. Differentiability: the right derivative is plainly $a$. For the left, the shortcut fails outright, because
\[
\lim_{x\to0^{-}}f'(x)=\lim_{x\to0^{-}}\left(2x\sin\frac1x-\cos\frac1x\right)
\]
does not exist. Go to the definition instead:
\[
\left|\frac{f(h)-f(0)}{h}\right|=\abs{h\sin(1/h)}\leq\abs{h}\To0 ,
\]
so $f'_{-}(0)=0$ and therefore $a=0$, $b=0$. The matched function is $f(x)=x^{2}\sin(1/x)$ patched by the zero function, differentiable at the junction with $f'$ discontinuous there.
\end{worked}

\begin{keynote}[Counting the conditions]
Matching to order $k$ at a junction imposes $k+1$ equations, one for each of $f,f',\dots,f^{(k)}$, and therefore needs $k+1$ free parameters to be solvable in general. If a problem hands you two parameters and the continuity equation turns out to be automatic, the intended condition is $C^{2}$, not $C^{1}$: count the equations you actually get, not the ones you expected.
\end{keynote}

\begin{trap}
Imposing differentiability without continuity. A function discontinuous at $c$ is never differentiable there, so a system built from the derivative condition alone can have a solution that solves nothing. The mirror error is matching derivatives by substituting $c$ into the two derivative \emph{formulas} when one of the pieces has no one-sided derivative at $c$: for $f(x)=x^{2}\sin(1/x)$ patched to $f(0)=0$, the formula for $f'$ has no limit at $0$ although $f'(0)$ exists.
\end{trap}

\begin{seealsob}One-sided derivatives, corners, cusps and vertical tangents; differentiable but not continuously differentiable; the derivative from first principles.\end{seealsob}

Piecewise definitions are often disguised. An absolute value is a two-case definition, a maximum of two functions is a two-case definition, and a fractional power is a formula whose derivative simply ceases to exist at one point. It is worth having a vocabulary for the three ways a derivative can fail while the function stays continuous, because problems ask you to name them.

\begin{keynote}[Four ways a continuous function fails to be differentiable]
\textbf{Corner:} both one-sided derivatives finite and unequal, as $\abs{x}$ at $0$; the graph has two distinct tangent directions. \textbf{Cusp:} both infinite with opposite signs, as $x^{2/3}$; the two branches leave vertically on the same side. \textbf{Vertical tangent:} both infinite with the same sign, as $x^{1/3}$; the graph is smooth as a curve but not as a function of $x$. \textbf{Oscillation:} neither one-sided limit exists, as $x\sin(1/x)$; there is no tangent direction at all. Continuity is assumed throughout, and a discontinuity is a fifth and cruder failure that needs no analysis.
\end{keynote}

\tech{One-sided derivatives: corners, cusps and vertical tangents}{medium}
\cue Absolute values, $\max$ or $\min$ of two functions, floors, or fractional powers: $\abs{x^{2}-4}$, $\min(x,x^{2})$, $x^{2/3}$, $\sqrt[3]{x}$.
\method Compute the two one-sided limits separately,
\[
f'_{-}(a)=\lim_{h\to0^{-}}\frac{f(a+h)-f(a)}{h},\qquad f'_{+}(a)=\lim_{h\to0^{+}}\frac{f(a+h)-f(a)}{h},
\]
and conclude that $f'(a)$ exists exactly when both exist finitely and agree. Then name the failure. \emph{Corner:} both finite, unequal, as $\abs{x}$ at $0$. \emph{Cusp:} both infinite with opposite signs, as $x^{2/3}$ at $0$. \emph{Vertical tangent:} both infinite with the same sign, as $x^{1/3}$ at $0$. For $\abs{g(x)}$ the failures sit exactly at the \emph{sign-changing} zeros of $g$, and at those points $f'_{\pm}=\pm g'$, so the jump is $2\abs{g'}$.

\begin{worked}
Let $f(x)=\abs{x^{2}-4}$. The inner function changes sign at $x=\pm2$ and nowhere else, so those are the only candidates. Away from them $f'(x)=2x\sgn(x^{2}-4)$. At $x=2$, with $g(x)=x^{2}-4$ and $g'(2)=4$:
\[
f'_{-}(2)=-g'(2)=-4,\qquad f'_{+}(2)=g'(2)=4 ,
\]
a corner with a jump of $8$. Similarly $f'_{-}(-2)=4$ and $f'_{+}(-2)=-4$.

Now the contrast that decides whether a sign change matters. Take $F(x)=\abs{x^{2}-4}\,(x-2)$ near $x=2$: the factor $(x-2)$ vanishes there, and
\[
\frac{F(2+h)-F(2)}{h}=\frac{\abs{(2+h)^{2}-4}\cdot h}{h}=\abs{4h+h^{2}}\To0 ,
\]
so $F'(2)=0$ exists despite the absolute-value bars. The bars are not the obstruction; an unmatched pair of one-sided slopes is.

And the three named failures, side by side at the origin. For $x^{2/3}$, $\dfrac{h^{2/3}}{h}=h^{-1/3}$, which tends to $+\infty$ from the right and $-\infty$ from the left: a cusp. For $x^{1/3}$, $\dfrac{h^{1/3}}{h}=h^{-2/3}\to+\infty$ from both sides: a vertical tangent, and the curve is smooth as a set even though $f'$ does not exist. For $\min(x,x^{2})$, the pieces cross at $x=0$ and $x=1$; at $0$ the slopes are $1$ and $2x=0$, a corner, and at $1$ they are $1$ and $2$, another corner.

One with a floor in it, to show that the method is unchanged. Let $f(x)=\floorb{x}+\bigl(x-\floorb{x}\bigr)^{2}$. On each interval $[n,n+1)$ this is $n+(x-n)^{2}$, which rises from $n$ to $n+1$, so $f$ is continuous at every integer despite the floor: the left limit at $n$ is $(n-1)+1=n$ and the value is $n$. Differentiability is another matter. On $(n-1,n)$ the derivative is $2(x-n+1)$, tending to $2$ as $x\to n^{-}$; on $(n,n+1)$ it is $2(x-n)$, tending to $0$ as $x\to n^{+}$. So $f'_{-}(n)=2$ and $f'_{+}(n)=0$: a corner at every integer, with $f$ continuous everywhere and differentiable off $\Z$.
\end{worked}

\begin{trap}
Declaring $\abs{x}^{3}$ non-differentiable at $0$ because of the bars. In fact $\abs{x}^{3}$ has derivative $3x\abs{x}$ and second derivative $6\abs{x}$, so it is $C^{2}$ at the origin and fails only at the third derivative. The kink comes from a sign change of the inside, never from the notation, and at a zero of even order there is no sign change. The reverse error is checking only the zeros of $g$ and forgetting that $\abs{g}$ can also fail where $g$ itself is not differentiable.
\end{trap}

\begin{seealsob}Piecewise matching for continuity and differentiability; differentiable but not continuously differentiable; sign charts for monotonicity and concavity.\end{seealsob}

One final distinction, and it is the one that most often goes unexamined. ``$f'(a)$ exists'' and ``$f'$ is continuous at $a$'' are different statements, and the second is strictly stronger. The standard family separating them is $x^{k}\sin(1/x)$, and it repays being worked through carefully once, because the same computation settles a whole row of examination questions about whether a limit of $f'$ says anything about $f'(a)$.

\tech[Differentiable but not continuously differentiable]{Differentiable but not $C^{1}$}{hard}
\cue A function of the form $x^{k}\sin(1/x)$ with $f(0)=0$, and a question about $f'(0)$, about $\lim_{x\to0}f'(x)$, or about whether $f$ is $C^{1}$.
\method Two separate computations, kept apart. First, $f'(0)$ from the definition, by squeezing the difference quotient. Second, $f'(x)$ for $x\neq0$ by the rules, followed by an examination of $\lim_{x\to0}f'(x)$. The exponent decides everything: $k=1$ gives no $f'(0)$, $k=2$ gives $f'(0)=0$ with $f'$ discontinuous at $0$, and each further power buys one more order of smoothness. Note also Darboux's theorem: a derivative has the intermediate value property, so $f'$ can never have a jump discontinuity, only an oscillatory one, which is exactly what this family provides.

\begin{worked}
Let $f(x)=x^{2}\sin(1/x)$ for $x\neq0$ and $f(0)=0$. From the definition,
\[
\left|\frac{f(h)-f(0)}{h}\right|=\abs{h\sin(1/h)}\leq\abs{h}\To0 ,\qquad\text{so } f'(0)=0 .
\]
Away from the origin the rules give
\[
f'(x)=2x\sin\frac1x-\cos\frac1x ,
\]
whose first term tends to $0$ but whose second oscillates between $-1$ and $1$ forever. So $\lim_{x\to0}f'(x)$ does not exist, while $f'(0)=0$ does. The function is differentiable on all of $\R$ and its derivative is discontinuous at one point.

The neighbouring cases. For $k=1$, the quotient is $\sin(1/h)$, which oscillates without settling, so $f'(0)$ does not exist at all; here $f$ is continuous but not differentiable. For $k=3$, $f'(x)=3x^{2}\sin(1/x)-x\cos(1/x)\to0=f'(0)$, so $f$ is $C^{1}$; but the second difference quotient $\frac{f'(h)-f'(0)}{h}=3h\sin(1/h)-\cos(1/h)$ has no limit, so $f''(0)$ does not exist. Each unit of $k$ buys one derivative, and the pattern continues.
\end{worked}

\begin{trap}
Concluding that $f'(0)$ fails to exist because $\lim_{x\to0}f'(x)$ fails to exist. Those are different statements and this family is the standard counterexample. The valid implication runs only one way: \emph{if} $f$ is continuous at $a$ and $\lim_{x\to a}f'(x)=L$ exists, \emph{then} $f'(a)=L$. The converse is false, and any argument that assumes it has assumed $C^{1}$ without saying so.
\end{trap}

\begin{seealsob}The derivative from first principles; one-sided derivatives, corners, cusps and vertical tangents; Rolle and mean-value arguments as tools.\end{seealsob}

\subsection{A protocol for differentiability questions}

Almost every question in this half of the chapter yields to the same four-step reading, and it is worth carrying explicitly, because the steps have to happen in this order.

\emph{First, is $f$ continuous at $a$?} If not, stop: no derivative exists and no matching condition can rescue it. This step is free and it eliminates a whole class of wrong answers.

\emph{Second, does a rule reach $a$?} A rule reaches $a$ when $f$ agrees with a single elementary formula on a whole neighbourhood of $a$. A junction point, an absolute value at a sign change, an assigned value filling a hole, and a fractional power at its branch point all fail this test, and each is an instruction to use the definition.

\emph{Third, compute both one-sided difference quotients.} Not the one-sided limits of $f'$: those agree with the one-sided derivatives only when they exist, and the whole point of the hard cases is that they may not. Where a factor can be bounded, squeeze rather than expand.

\emph{Fourth, say precisely what was found.} ``$f'(a)$ exists and equals $L$'' is a different statement from ``$f'$ is continuous at $a$'', which is different again from ``$\lim_{x\to a}f'(x)=L$''. A question that mentions $C^{1}$, or asks whether $f'$ is continuous, is asking the second; a question that asks for $f'(a)$ is asking the first. Answering the wrong one is the commonest way to lose a mark while doing correct mathematics.

Two habits carry this chapter. The first is counting the doors: whenever $x$ appears in an integral, ask through how many of the three entrances it enters, and produce exactly that many terms. The second is knowing when to stop looking for a rule. A junction point, an assigned value at a removable singularity, an absolute value at a sign change, and a deliberately oscillatory patch are all announcements that the rule systems have run out, and each one is answered by the same difference quotient, computed one side at a time.

\section{Recognition matrix}
\begin{recog}{@{}p{0.30\textwidth}p{0.36\textwidth}p{0.28\textwidth}@{}}
\rowcolor{mist}\mh{If you see} & \mh{Reach for} & \mh{If that fails} \\
\midrule
\endhead
$\int_{a}^{u(x)}f(t)\dd t$ & $f(u(x))u'(x)$; never find the antiderivative & Split at a constant if both limits move \\
$\int_{u(x)}^{a}f(t)\dd t$ & Flip orientation first: $-f(u(x))u'(x)$ & Write $b'$ and $a'$ with signs before substituting \\
$\int_{a(x)}^{b(x)}f(t)\dd t$ & $f(b)b'-f(a)a'$ & Look for a complementary identity in the two terms \\
Reciprocal limits $\int_{1/x}^{x}$ & Expect $\arctan x+\arctan(1/x)=\pi/2$ or a sibling & Differentiate and simplify before evaluating \\
$x$ in the limits \emph{and} in the integrand & Full Leibniz: two boundary terms plus $\int\pdv{F}{x}$ & Count the doors; three doors, three terms \\
The same $\lambda(x)$ in a limit and multiplying $t$ & Substitute $s=\lambda(x)t$; carry the Jacobian & Fall back on the three-term rule \\
Fixed limits, a parameter in the integrand & Differentiate under the sign; boundary terms vanish & Recognise a known parametric identity first \\
An integral you cannot evaluate, with a parameter & Get $I'(a)$, integrate back, fix $C$ at an easy $a$ & Check uniform convergence before swapping \\
$\int_{0}^{x}K(x-t)y(t)\dd t=h(x)$ & Differentiate if $K$ is regular; Laplace if $K$ is singular & $\widehat{u^{-1/2}}=\sqrt{\pi/s}$ for the Abel kernel \\
``From the definition'' & $\lim_{h\to0}\frac{f(a+h)-f(a)}{h}$, honestly & Squeeze: $\abs{\Delta f}\leq C\abs{h}^{1+\epsilon}$ gives $f'(a)=0$ \\
$\lim\frac{g(a+\alpha h)-g(a+\beta h)}{h}$ & $(\alpha-\beta)g'(a)$ & Identify $g$ by the argument that repeats \\
$n\bigl(g(a+c/n)-g(a)\bigr)$ as $n\to\infty$ & Same limit with $h=1/n$: answer $cg'(a)$ & Not L'Hopital; it needs $g'$ near $a$ \\
$f(a)$ assigned to fill a $1^{\infty}$ hole & Expand $\ln f$ to second order; $f'(a)=f(a)c_{1}$ & First order is one order too few \\
Two formulas meeting at $c$, unknown $a,b$ & Continuity first, then $f'_{-}(c)=f'_{+}(c)$ & Matching $k$ derivatives needs $k+1$ equations \\
$\abs{g(x)}$ & Non-differentiable only at sign-changing zeros of $g$ & Jump in slope is $2\abs{g'}$ there \\
$x^{2/3}$, $x^{1/3}$ at $0$ & Cusp: $\pm\infty$ opposite; vertical tangent: same sign & Both are one-sided limits, computed separately \\
$x^{k}\sin(1/x)$ with $f(0)=0$ & Definition for $f'(0)$; rules for $f'(x)$, $x\neq0$ & $k=1$ none, $k=2$ exists but $f'$ jumps \\
$\lim_{x\to a}f'(x)$ does not exist & Says nothing about $f'(a)$ & Darboux: a derivative has no jump discontinuity \\
$G(x)=\int_{a}^{x}f$ with $f>0$, inverse wanted & $\bigl(G^{-1}\bigr)'(y)=1/f(x)$ at the $x$ with $G(x)=y$ & Locating $x$ is root-finding, not calculus \\
Concavity of $\int_{a}^{u(x)}f$ & $F''=f'(u)u'^{2}+f(u)u''$ & Two terms; the second dies only if $u$ is affine \\
A floor or fractional part in a formula & Rewrite as cases on $[n,n+1)$, then one-sided limits & Continuity at the joins is not automatic \\
\end{recog}
